\documentclass[../../main.tex]{subfiles}
\begin{document}

\begin{center}
  {\fontsize{48pt}{58pt}\selectfont \textbf{2010}\normalsize\textbf{年度}}
\end{center}

\begin{table}[htb]
  \centering
  \caption{出題テーマと難易度}
  \begin{tabular}{|c|p{9cm}|c|}
    \hline
    \textbf{問題番号} & \textbf{テーマ} & \textbf{難易度} \\
    \hline
    第1問 & 方程式の解の存在証明と定積分の評価 & \\
    \hline
    第2問 & ガウス記号を含む方程式と無限級数 & \\
    \hline
    第3問 & 確率（3の倍数条件のカード） & \\
    \hline
    第4問 & 領域の図示（ベクトル条件と円） & \\
    \hline
  \end{tabular}
\end{table}

\begin{table}[htb]
  \centering
  \caption{解答}
  \renewcommand{\arraystretch}{1.4}
  \begin{tabular}{|c|c|p{9cm}|}
    \hline
    \textbf{問題番号} & \textbf{小問} & \textbf{答え} \\
    \hline
    第1問 & (1) & 証明問題 \\
    \cline{2-3}
    & (2) & $J=2\alpha^2\sin\alpha-4\alpha+4\sin\alpha$ \\
    \cline{2-3}
    & (3) & 証明問題（$J>\sqrt{2}$） \\
    \hline
    第2問 & (1) & $a=7$：$x=2$，$a=8$：解なし，$a=9$：$x=3$ \\
    \cline{2-3}
    & (2) & $a_1=3$，$a_2=8$ \\
    \cline{2-3}
    & (3) & $\dfrac{3}{4}$ \\
    \hline
    第3問 & (1) & $p(8)=\dfrac{1}{4}$ \\
    \cline{2-3}
    & (2) & $p(3k+2)=\dfrac{k}{3k+2}$ \\
    \hline
    第4問 & - & $D=\left\{(x,y)\ \middle|\ x\ge0,\ \left(x-\dfrac{a}{2}\right)^2+y^2\ge\left(\dfrac{a}{2}\right)^2\right\}\setminus\{(0,0),(a,0)\}$ \\
    \hline
  \end{tabular}
\end{table}

\paragraph{試験形態}
4問．試験時間は要確認．

\end{document}
